\begin{conclusion}[坏素数 \(37\) 处的 Latt\`es 稳定约化：公共平方因子、Chebyshev 共轭与结点退化的同步]\label{con:xi-terminal-zm-elliptic-lattes-badprime-37-degree-drop}
沿用结论 \ref{con:xi-terminal-zm-elliptic-lattes-doubling} 的倍点 Latt\`es 自映射
\[
\Phi(X)=\frac{X^4+2X^2-2X+1}{4X^3-4X+1}\in\QQ(X).
\]
则在 \(\FF_{37}[X]\) 中出现非平凡公共因子，且满足显式同余分解
\[
X^4+2X^2-2X+1\equiv (X-5)^2\,(X^2+10X+3)\pmod{37},
\]
\[
4X^3-4X+1\equiv (X-5)^2\,(4X+3)\pmod{37}.
\]
因此在稳定约化意义下（约去公共因子）得到次数严格下降的约化映射
\[
\boxed{\;
\overline{\Phi}(X)=\frac{X^2+10X+3}{4X+3}\in\FF_{37}(X),\qquad \deg\overline{\Phi}=2<\deg\Phi=4\; } .
\]
\par
并且在 \(\PP^{1}/\FF_{37}\) 上存在显式 Möbius 变换
\[
M(X)=\frac{2X+13}{X-5}\in \mathrm{PGL}_{2}(\FF_{37})
\]
使得
\[
\boxed{\;
M\circ \overline{\Phi}\circ M^{-1}(X)=X^{2}-2\; } .
\]
因而 \(\overline{\Phi}\) 在 \(\FF_{37}\) 上与归一化 Chebyshev 多项式 \(T_{2}(X)=X^{2}-2\) 共轭。
在该约化动力系统中，\(X=5\) 满足 \(\overline{\Phi}(5)=5\) 且 \(\overline{\Phi}'(5)=0\)，从而为唯一的临界不动点；
同时 \(M\) 在 \(X=5\) 处具有极点并将其送至 \(\infty\)，从而把 \(\overline{\Phi}\) 刚性化为多项式正规形。
\par
另一方面，Lee--Yang 分歧点的三次控制多项式（结论 \ref{con:xi-terminal-zm-leyang-ramification-critical-values}）
\[
g(X):=16X^3-9X^2+1
\]
在 \(\FF_{37}[X]\) 中发生双根塌缩：
\[
g(X)\equiv 16\,(X-16)(X-5)^2\pmod{37}.
\]
并且 Lee--Yang 三次因子在 \(\FF_{37}[y]\) 上同样出现重根融合（结论 \ref{con:xi-terminal-zm-leyang-badprime-37-common-double-root}）：
\[
P_{\mathrm{LY}}(y)\equiv -3\,(y-14)(y-6)^2\pmod{37}.
\]
\par
因此判别式二次脊曲线
\[
C:\quad W^{2}=-y(y-1)P_{\mathrm{LY}}(y)
\]
在 \(p=37\) 处出现由二重根 \(y\equiv 6\) 所触发的半稳定退化：
其特殊纤维满足
\[
C_{\FF_{37}}:\quad W^{2}=3\,y(y-1)(y-14)(y-6)^{2},
\]
为不可约结点曲线。
令 \(Z=W/(y-6)\)，则归一化为椭圆曲线
\[
E_{14}:\quad Z^{2}=3\,y(y-1)(y-14)
\]
并且该结点为非分裂结点（两支仅在 \(\FF_{37^{2}}\) 上分离）。
从而对任意 \(n\ge 1\) 有点数恒等式
\[
\card{C(\FF_{37^{n}})}=\card{E_{14}(\FF_{37^{n}})}+(-1)^{n+1},
\]
并得到局部 \(\zeta\)-函数的显式分解
\[
Z_{C,37}(T)=(1+T)\,Z_{E_{14},37}(T),\qquad
Z_{E_{14},37}(T)=\frac{1-2T+37T^{2}}{(1-T)(1-37T)}.
\]
\par
综上，素数 \(37\) 同时承载：倍点 Latt\`es 自映射的稳定约化降阶与 Chebyshev 共轭、Lee--Yang 三次分歧的重根碰撞、以及判别式二次脊曲线的非分裂结点退化；
这些同步现象共同刻画了该椭圆化机制在唯一坏素处的算术奇异性。
\end{conclusion}
\begin{proof}[证明要点]
模 \(37\) 的因式分解与约去为有限整系数计算。
Möbius 共轭可在 \(\FF_{37}(X)\) 中直接代入化简核验。
结点曲线的归一化由代换 \(Z=W/(y-6)\) 给出，而非分裂性等价于归一化纤维在 \(y=6\) 处的两点仅在二次扩张上可见，从而点数差异呈 \((-1)^{n+1}\) 交替并推出所示 \(\zeta\)-因子。
可复现核验脚本：\texttt{scripts/exp\_fold\_zm\_elliptic\_lattes\_rational\_points\_audit.py}（\(\overline{\Phi}\) 的约化、Chebyshev 共轭与周期枚举），\texttt{scripts/exp\_fold\_zm\_elliptic\_leyang\_cover\_geometry\_audit.py}（\(E_{14}\) 点数与 \((-1)^{n+1}\) 修正）。
最后一条重根分解由结论 \ref{con:xi-terminal-zm-leyang-badprime-37-common-double-root} 给出。证毕。
\end{proof}

\begin{conclusion}[三除点重量像在 \texorpdfstring{$\FF_{37}$}{F37} 上的六重坍缩]\label{con:xi-terminal-zm-elliptic-3torsion-y-image-mod37-collapse}
仍沿用结论 \ref{con:xi-terminal-zm-elliptic-3torsion-y-image-octic} 的八次消元多项式
\[
L_{3}(y)=81 y^{8}-324 y^{7}+297 y^{6}+729 y^{5}-216 y^{4}-1395 y^{3}-480 y^{2}-120 y+7\in\ZZ[y].
\]
则在 \(\FF_{37}[y]\) 中发生高重塌缩分解
\[
L_{3}(y)\equiv 7\,(y-6)^{6}\,(y^{2}-5y-1)\pmod{37}.
\]
特别地，\(y\equiv 6\) 为 \(y\bigl(E[3]\setminus\{O\}\bigr)\) 的模 \(37\) 主聚集像，其重数为 \(6=3^{2}-3\)。
\end{conclusion}
\begin{proof}[证明要点]
对 \(L_{3}(y)\) 作模 \(37\) 约化并在 \(\FF_{37}[y]\) 中因式分解即可得到所示同余分解。
可复现核验脚本：\texttt{scripts/exp\_fold\_zm\_elliptic\_odd\_torsion\_mod37\_collapse\_audit.py}。证毕。
\end{proof}

\begin{conclusion}[五除点重量像的次数 \(24\) 消元证书与模 \(37\) 的二十重坍缩]\label{con:xi-terminal-zm-elliptic-5torsion-y-image-degree24-mod37-collapse}
仍在椭圆曲线
\(
E:\ Y^{2}=X^{3}-X+\tfrac14
\)
与重量函数
\(
y=X^{2}+Y-\tfrac12
\)
的设定下，取 \(5\)-division 多项式（\(X\)-坐标形式）
\[
\psi_{5}(X)=5X^{12}-62X^{10}+95X^{9}-105X^{8}-60X^{7}+285X^{6}-174X^{5}-5X^{4}-5X^{3}+35X^{2}-15X+2\in\ZZ[X].
\]
令
\[
L_{5}(y):=\operatorname{Res}_{X}\bigl(\psi_{5}(X),\Pi(X,y)\bigr)\in\ZZ[y],
\qquad
\Pi(X,y):=X^{4}-X^{3}-(2y+1)X^{2}+X+y(y+1).
\]
则 \(L_{5}(y)\) 在 \(\QQ\) 上不可约，且 \(\deg L_{5}=24\)；并且可取其闭式证书为
\[
\boxed{\;
\begin{aligned}
L_{5}(y)=\;&625 y^{24}-23500 y^{23}+244475 y^{22}+855915 y^{21}-11037859 y^{20}-89302755 y^{19}-22555640 y^{18}\\
&+761997855 y^{17}+1386419000 y^{16}-320647060 y^{15}-1928743150 y^{14}+838813180 y^{13}+3910018680 y^{12}\\
&+3220925750 y^{11}+1313468808 y^{10}-538699445 y^{9}-1045141370 y^{8}-317786725 y^{7}-89515935 y^{6}\\
&-310993491 y^{5}-240190705 y^{4}-86874340 y^{3}-15726875 y^{2}-1254110 y+10048.
\end{aligned}
\;}
\]
在坏素数 \(37\) 处，其模约化发生指数为 \(20=5^{2}-5\) 的集中塌缩：
\[
L_{5}(y)\equiv -4\,(y-6)^{20}\,\bigl(y^{4}+y^{3}+10y^{2}+9y+4\bigr)\pmod{37}.
\]
\end{conclusion}
\begin{proof}[证明要点]
由结果式的基本性质，\(L_{5}(y)\) 的根集合恰为 \(y\bigl(E[5]\setminus\{O\}\bigr)\) 的消元像。
不可约性与同余分解可由有限代数计算核验；见脚本 \texttt{scripts/exp\_fold\_zm\_elliptic\_odd\_torsion\_mod37\_collapse\_audit.py} 的可复现输出。证毕。
\end{proof}

\begin{conclusion}[七除点重量像在模 \(37\) 下的四十二重坍缩]\label{con:xi-terminal-zm-elliptic-7torsion-y-image-mod37-collapse}
令 \(\psi_{7}(X)\in\ZZ[X]\) 为同一曲线 \(E\) 的 \(7\)-division 多项式（\(X\)-坐标形式，次数 \(24\)）。
在有限域 \(\FF_{37}\) 上定义
\[
L_{7}^{(37)}(y):=\operatorname{Res}_{X}\bigl(\psi_{7}(X),\Pi(X,y)\bigr)\in\FF_{37}[y],
\]
其中 \(\Pi(X,y)\) 如上。
则在 \(\FF_{37}[y]\) 中存在精确分解
\[
L_{7}^{(37)}(y)= -4\,(y-6)^{42}\,\bigl(y^{6}-11y^{5}+18y^{4}+7y^{3}-6y^{2}-13y-9\bigr),
\]
其中指数 \(42=7^{2}-7\)，剩余因子次数 \(6=7-1\)。
\end{conclusion}
\begin{proof}[证明要点]
在 \(\FF_{37}[X,y]\) 中计算结果式并因式分解即可得到所示恒等式；见脚本 \texttt{scripts/exp\_fold\_zm\_elliptic\_odd\_torsion\_mod37\_collapse\_audit.py}。证毕。
\end{proof}

\begin{conjecture}[非分裂乘法约化素点处的普适坍缩指数]\label{conj:xi-terminal-zm-elliptic-odd-torsion-mod37-collapse-exponent}
对任意奇整数 \(n\ge 3\)（与 \(37\) 互素），令 \(\psi_{n}(X)\) 为曲线 \(E\) 的 \(n\)-division 多项式（\(X\)-坐标形式），并定义
\[
L_{n}(y):=\operatorname{Res}_{X}\bigl(\psi_{n}(X),\Pi(X,y)\bigr)\in\ZZ[y].
\]
则在 \(\FF_{37}[y]\) 中存在单位 \(c_{n}\in\FF_{37}^{\times}\) 与次数 \(n-1\) 的多项式 \(R_{n}(y)\) 使得
\[
L_{n}(y)\equiv c_{n}\,(y-6)^{n^{2}-n}\,R_{n}(y)\pmod{37},
\qquad
\deg R_{n}=n-1,
\qquad
R_{n}(6)\neq 0.
\]
\end{conjecture}
\begin{remark}
若把 \(p=37\) 处的非分裂乘法约化通过 Tate 曲线统一化理解，则在 \(\overline{\FF}_{37}\) 上可将 \(E[n]\) 的几何成分分解为 \(\mu_{n}\times (\ZZ/n\ZZ)\) 的两类参数族；其中后者在约化时趋向同一奇点分量，从而在某一重量像值（此处为 \(y\equiv 6\)）产生 \(n(n-1)=n^{2}-n\) 阶的集中坍缩，而 \(\mu_{n}\setminus\{1\}\) 的 \(n-1\) 个点贡献剩余次数 \(n-1\) 的因子。
该机制与上文 \(P_{\mathrm{LY}}(y)\) 在模 \(37\) 下于 \(y=6\) 出现二重根的退化现象一致。
\end{remark}

